\documentclass[UTF8,a4paper,12pt]{ctexrep}

\usepackage{geometry}
\usepackage{amsmath}
\usepackage{amssymb}
\usepackage{booktabs}
\usepackage{longtable}
\usepackage{array}
\usepackage{graphicx}
\usepackage{hyperref}
\usepackage{caption}
\usepackage{subcaption}
\usepackage{enumitem}
\usepackage{siunitx}
\usepackage[numbers,sort&compress]{natbib}

\geometry{left=3cm,right=2.5cm,top=2.5cm,bottom=2.5cm}
\hypersetup{
  colorlinks=true,
  linkcolor=black,
  citecolor=black,
  urlcolor=blue
}

\sisetup{
  detect-all=true,
  per-mode=symbol
}

\title{面向数据中心冷却系统的 TES-PV-TOU MILP-MPC 控制方法研究}
\author{作者姓名}
\date{\today}

\begin{document}

\maketitle

\begin{abstract}
数据中心具有连续运行、热负荷密度高和冷却可靠性要求高等特点，冷却系统运行策略直接影响数据中心的能耗、运行成本和温度安全。随着分布式光伏和分时电价信号的引入，数据中心冷却系统不再只是被动维持机房温度，而是需要在未来负荷、天气、可再生能源出力和电价变化下进行动态优化调度。蓄冷系统能够在时间上解耦制冷功率和冷量需求，为削峰填谷、提高光伏消纳和降低运行成本提供了物理基础。

本文面向南京数据中心冷却场景，构建包含冷机模式、冷水蓄冷系统、光伏出力、分时电价和机房温度约束的控制导向模型，并设计基于混合整数线性规划的模型预测控制方法。该方法在 12 h 滚动预测时域内优化冷机出力、TES 充放冷功率、全设施侧电网购电功率和 PV 弃电功率，同时显式考虑冷机 mode 约束、TES SOC 约束、signed valve 充放冷约束、温度安全约束和终端状态要求。当前可复现 7 day 归因结果表明，direct no-TES、price-aware no-TES MPC、rule-based TES 和 TES-MPC 均可稳定运行，且 fallback、温度越界、物理一致性违约和 signed-valve 违约均为零。在全设施侧主成本口径下，price-aware no-TES MPC 相比 direct no-TES 降本约 4.56\%，而 TES-MPC 相比 no-TES MPC 的增量降本约 0.09\%。新增 30 day 中国 TOU/DR 矩阵显示，在 61 组配对场景中 TES-MPC 相比 no-TES MPC 的月平均增量节省为 518.81 CNY，bootstrap 置信区间为 313.84 到 737.38 CNY。三地点 PV--TES 技术容量筛选已使用南京、广州和北京真实 EPW、PVGIS 20 MWp 光伏曲线和江苏 TOU 电价完成 EnergyPlus-derived 年度矩阵；但该矩阵只能作为筛选基线，不能写作最终 MPC+EnergyPlus 在线耦合主结果。当前修正版已完成在线 MPC+EnergyPlus smoke/pilot，验证 MPC 动作可通过 Runtime API 写入 EnergyPlus 并回读温度、TES 和全设施电功率响应；完整三地点年度在线矩阵仍待运行。因此，本文将大部分经济收益归因于价格感知冷机调度、热惯性和 mode 优化，并将 TES 的增量贡献、容量边界及在线耦合验证状态单独报告。

\textbf{关键词：} 数据中心冷却；蓄冷系统；光伏消纳；分时电价；模型预测控制；混合整数线性规划
\end{abstract}

\chapter*{Abstract}
Data centers operate continuously and require reliable cooling to maintain safe server-room temperatures. With the integration of photovoltaic generation and time-of-use electricity tariffs, the cooling system becomes a dynamic dispatch problem rather than a simple temperature-regulation problem. Thermal energy storage can decouple cooling production from cooling demand over time, which creates opportunities for peak shaving, renewable-energy utilization, and operating-cost reduction.

This thesis develops a control-oriented framework for a data-center cooling system with chiller plant modes, chilled-water thermal energy storage, photovoltaic generation, and time-of-use electricity prices. A mixed-integer linear model predictive controller is formulated to optimize chiller load, TES charging and discharging, whole-facility grid import, and PV curtailment over a 12 h receding horizon while enforcing temperature, SOC, signed-valve, and mode-switching constraints. Reproducible 7 day attribution results show that direct no-TES, price-aware no-TES MPC, rule-based TES, and TES-MPC cases all run with zero fallback, zero temperature degree-hour violation, and zero physical or signed-valve violation. Under the whole-facility cost boundary, price-aware no-TES MPC reduces total cost by about 4.56\% relative to direct no-TES, while TES-MPC provides an incremental cost reduction of about 0.09\% relative to no-TES MPC. A new 30 day China TOU/DR matrix shows an average monthly incremental saving of 518.81 CNY for TES-MPC relative to no-TES MPC across 61 paired scenarios, with a bootstrap confidence interval from 313.84 to 737.38 CNY. A three-location PV--TES technical screening matrix has been completed with real EPW weather, PVGIS 20 MWp PV profiles, and the Jiangsu TOU tariff, but that result is an EnergyPlus-derived profile replay baseline rather than the final online MPC+EnergyPlus result. The corrected online MPC+EnergyPlus path has been verified with smoke and pilot runs: MPC actions are written through the EnergyPlus Runtime API and physical temperature, TES, and whole-facility power responses are read back from EnergyPlus. The full-year three-location online sizing matrix remains to be run. The results therefore separate economic chiller scheduling from the incremental contribution, sizing boundary, and online-coupling validation status of TES.

\textbf{Keywords:} data center cooling; thermal energy storage; photovoltaic self-consumption; time-of-use tariff; model predictive control; mixed-integer linear programming

\tableofcontents

\chapter{绪论}

\section{研究背景}

数据中心是数字基础设施的重要组成部分，其 IT 设备通常需要全年连续运行。服务器、存储和网络设备产生的热量需要由冷却系统及时移除，否则机房温度升高会影响设备可靠性甚至造成运行风险。因此，数据中心冷却系统不仅承担能耗管理任务，也承担安全约束管理任务。

与普通办公建筑空调相比，数据中心冷却负荷更加稳定，温度约束更加刚性，运行中断代价更高。冷却系统的运行策略如果只追求短时温度控制，可能忽略未来电价、光伏出力和蓄冷状态变化；如果只追求经济性，又可能造成温度或设备约束风险。因此，该问题天然适合被表述为带约束的动态优化控制问题。

近年来，分布式光伏和分时电价逐渐进入数据中心能源管理场景。光伏出力具有明显的日内波动和预测误差，分时电价则引导用户将可转移负荷从高价时段转移到低价时段。对于冷却系统而言，冷水蓄冷系统可以把低价或高光伏时段产生的冷量存储起来，在高价或光伏不足时段释放，从而使冷却能耗具有一定的时间可调性。

\section{数据中心冷却系统能耗问题}

数据中心总能耗通常由 IT 负荷、冷却系统、供配电损耗和辅助设备能耗构成。PUE 是评价数据中心能源效率的常见指标，其基本定义为总设施功率与 IT 功率之比：

\begin{equation}
  PUE_k = \frac{P_{facility,k}}{P_{IT,k}}.
\end{equation}

当 IT 负荷给定时，降低冷却功率和辅助设备功率可以降低 PUE。然而，在包含 TES、PV 和 TOU 电价的系统中，单纯降低某一时刻的冷却功率并不一定等价于降低总运行成本。系统可能在低价时段增加制冷功率为 TES 充冷，在高价时段降低电网购电功率。因而，论文需要同时评价运行成本、PV 消纳、温度约束和 TES 行为，而不能只评价单点能效。

\section{蓄冷系统、光伏与分时电价的调度价值}

TES 的核心价值在于冷量的时间转移。对于冷水蓄冷系统，充冷过程通常对应冷机额外运行并提升 TES SOC；放冷过程则利用已储存冷量满足部分冷却需求，降低当前冷机或电网用电压力。在 TOU 电价下，TES 可用于低价充冷、高价放冷；在 PV 场景下，TES 还可吸收白天富余光伏，减少弃光。

但 TES 并非在所有条件下都一定带来经济收益。充放冷效率损失、冷机 COP、PV 规模、TOU 价差、温度约束、终端 SOC 要求和 TES 容量功率配置都会影响最终收益。因此，本研究不预设“有 TES 必然省钱”，而是通过闭环仿真分析 TES-MPC 在不同条件下的行为和边界。

\section{传统规则控制的局限性}

传统蓄冷控制通常基于固定规则，例如夜间低价充冷、白天高价放冷、SOC 低于阈值后停止放冷、机房温度接近上限时增强制冷。这类规则易于实现，但难以同时处理未来 PV 出力、电价、天气、IT 负荷、SOC 约束、温度约束和 PV 弃电问题。

MPC 的优势在于利用未来预测，在滚动时域内优化控制动作，并显式处理系统约束。对于数据中心冷却系统，这种约束可解释性尤其重要，因为温度安全和设备运行边界不能依赖大量试错来学习。已有研究已将 MPC 用于大型冷站和 TES 调度\citep{ma2009tesMpc,ma2010buildingCoolingMpc,kim2022mpcTesPv}，也有研究讨论数据中心冷却 MPC 与冷水蓄冷控制\citep{lazic2018datacenterMpc,zhu2024hybridCoolingStorage,zhao2024multiChillerMpc}。本文在此基础上将 TES、PV、TOU 电价和温度约束统一到一个面向数据中心场景的控制框架中。

\section{本文研究目标与研究内容}

本文的研究目标是：面向南京数据中心冷却系统，建立包含 TES、PV、TOU 电价和机房温度约束的控制导向模型，设计基于 MILP 的 MPC 控制器，并通过多场景闭环仿真评估该方法在运行成本、PV 消纳、PUE、温度约束和鲁棒性方面的表现。

本文研究内容包括：

\begin{enumerate}[label=(\arabic*)]
  \item 建立数据中心冷却蓄冷系统边界，定义外部输入、状态变量、控制变量和评价指标。
  \item 构建 TES SOC 动态模型、机房温度代理模型、PV-grid 功率平衡模型和运行成本模型。
  \item 将 TES 充放冷互斥、SOC 限制、温度限制和 PV 弃电纳入 MILP-MPC 优化问题。
  \item 设计 no-TES baseline、rule-based TES 和 MILP-MPC 对比实验。
  \item 通过 PV 预测误差、hot/mild week 和 tariff sensitivity 场景分析方法的鲁棒性和适用边界。
  \item 使用真实 EPW、PVGIS 光伏曲线和江苏 TOU 电价构建 PV--TES 容量矩阵，给出运行层面的 technical recommended capacity range。
\end{enumerate}

\section{本文主要贡献}

本文的主要贡献预期包括：

\begin{enumerate}[label=(\arabic*)]
  \item 建立了面向数据中心冷却系统的 TES-PV-TOU 控制导向模型，明确了系统边界、关键状态和外部输入。
  \item 构建了考虑温度约束、SOC 约束、PV 弃电和分时电价的 MILP-MPC 控制器。
  \item 通过 baseline 对比和多场景闭环仿真，分析了 TES-MPC 控制策略的经济性、鲁棒性和适用边界。
\end{enumerate}

\section{论文结构安排}

第 1 章介绍研究背景、问题动机和本文贡献。第 2 章综述数据中心冷却、TES 技术和 MPC 控制相关研究。第 3 章建立数据中心冷却蓄冷系统的控制导向模型。第 4 章给出 MILP-MPC 控制方法。第 5 章设计实验场景和评价指标。第 6 章给出当前闭环仿真结果并讨论其边界。第 7 章总结研究结论、局限性和后续工作。

\chapter{文献综述}

\section{数据中心冷却系统与能耗评价}

数据中心冷却系统的首要目标是维持设备运行所需的温度和湿度环境。与舒适性空调相比，数据中心冷却更强调连续性、可靠性和故障容忍能力。数据中心能耗评价常使用 PUE、冷却系统能耗占比、峰值功率、温度越界次数等指标。近年来，综述研究也开始从优化控制、评价指标和工程可实施性角度系统总结数据中心冷却控制方法\citep{chang2024datacenterCoolingReview}。

在论文写作中，本节应重点说明数据中心冷却不是普通建筑空调问题。它具备高负荷密度、连续运行和安全约束强的特点，因此控制策略需要同时关注经济性和约束可行性。

\section{数据中心蓄冷技术研究现状}

数据中心 TES 技术包括冷水蓄冷、冰蓄冷和相变材料蓄冷等形式。相关综述指出，TES 可用于缓解冷量供需在时间、温度或空间上的不匹配，并可为数据中心提供削峰填谷、备用冷源和可再生能源消纳能力\citep{liu2020tesDatacenterReview}。

本文聚焦冷水蓄冷系统。该选择的原因是冷水蓄冷更容易与冷站系统和控制导向模型衔接，且可以用 SOC 动态、最大充放冷功率和效率损失等变量形成相对清晰的优化模型。其他 TES 技术在绪论和综述中简要说明即可，不应分散主线。

\section{建筑冷站系统中的模型预测控制}

建筑冷站和大型校园冷却系统中的 TES-MPC 研究为本文提供了直接方法基础。Ma 等研究了带冷水蓄冷罐的建筑冷却系统 MPC，证明控制导向模型与滚动优化可以用于冷站运行调度\citep{ma2009tesMpc,ma2010buildingCoolingMpc}。Kim 等进一步在大型冷站中将 TES、PV 和电网减碳目标结合，以 MILP-MPC 形式开展现场示范和 baseline 对比\citep{kim2022mpcTesPv}。

这些研究的共同特点是：建立简化但可解释的系统模型，显式考虑设备约束和储能状态，并通过 baseline comparison 评价 MPC 的实际增量。本文采用类似叙事逻辑，但将应用对象明确为数据中心冷却系统。

\section{数据中心冷却系统中的模型预测控制}

数据中心冷却 MPC 已被用于服务器机房温度、气流和冷站能耗控制。Lazic 等指出，真实物理系统中的强化学习部署会受到 unexpected events、limited data 和 expensive failures 等问题影响，因此采用 data-driven model-based/MPC 思路控制数据中心冷却\citep{lazic2018datacenterMpc}。Zhao 等研究了考虑整系统节能的数据中心多冷机 MPC\citep{zhao2024multiChillerMpc}。Zhu 等提出了数据中心冷水蓄冷混合冷却系统的先进控制策略，并使用 MILP 优化不同冷却模式\citep{zhu2024hybridCoolingStorage}。

这些文献说明数据中心冷却系统具备用 MPC 进行动态控制的研究基础。本文不采用 DRL 作为主线，是因为数据中心冷却具有强安全约束和高试错成本；MPC 在可解释性、约束处理和工程验证方面更适合当前毕业设计目标。

\section{蓄冷、光伏与分时电价协同调度研究}

TES、PV 和 TOU 电价的协同调度本质上是将能源供给、储能状态和用电价格统一到一个动态优化问题中。PV 出力高时，系统可选择增加充冷以吸收本地可再生能源；电价高时，系统可利用 TES 放冷减少电网购电；预测窗口内若未来高温或高价时段将出现，控制器应提前保留或补充 SOC。

已有大型冷站研究已将 PV 和 TES 纳入 MPC 框架\citep{kim2022mpcTesPv}。但在数据中心场景下，还需要更突出温度安全、IT 负荷连续性和 PUE 指标。本文的定位是将大型冷站 TES-PV-MPC 的建模和验证逻辑迁移到数据中心冷却蓄冷场景，并围绕可复现闭环仿真展开。

\section{文献总结与本文研究定位}

综上，已有研究分别覆盖了数据中心冷却控制、数据中心 TES 技术、建筑冷站 TES-MPC，以及 PV/电价下的冷站优化。但是，针对数据中心冷却系统，将 TES、PV、TOU 电价、机房温度约束和 SOC 约束统一到一个可复现实验框架中，并通过多场景闭环仿真分析经济性和适用边界，仍有进一步展开空间。

本文的研究定位不是提出全新的优化算法，而是建立一个边界清楚、公式完整、实验可复现的 TES-PV-TOU MILP-MPC 控制框架，并用该框架解释数据中心冷却蓄冷调度的机制和边界。

\chapter{数据中心冷却蓄冷系统建模}

\section{系统边界与总体结构}

本文研究对象为数据中心冷却蓄冷系统的经济型预测控制。系统包括 IT 负荷、机房温度、冷却系统、冷水 TES、PV、电网和 TOU 电价。外部扰动包括 IT 负荷、室外温度、PV 出力和电价；控制器通过 TES 充放冷和购电决策影响系统运行。

系统结构可以抽象为：

\begin{center}
\begin{tabular}{c}
Outdoor Temperature, IT Load, PV, TOU Price \\
$\downarrow$ \\
Cooling Demand and Energy Dispatch \\
$\downarrow$ \\
Chiller / TES / Grid Interaction \\
$\downarrow$ \\
Room Temperature, SOC, Cost, PV Utilization
\end{tabular}
\end{center}

当前仓库中的 \texttt{Nanjing-DataCenter-TES-EnergyPlus} 已提供 EnergyPlus 最小模型包、南京/北京/广州 EPW 天气文件、TOU 电价输入和 PVGIS 光伏输入。本文控制模型以控制导向代理为核心；最新代码已提供 EnergyPlus Runtime API 在线耦合 runner，但完整三地点年度在线矩阵仍待运行。

\section{外部输入变量}

\begin{table}[htbp]
  \centering
  \caption{外部输入变量}
  \begin{tabular}{llll}
    \toprule
    符号 & 含义 & 单位 & 说明 \\
    \midrule
    $P_{IT,k}$ & IT 负荷 & kW & 数据中心计算负荷 \\
    $T_{out,k}$ & 室外温度 & \si{\degreeCelsius} & 天气扰动 \\
    $P_{PV,k}$ & 光伏出力 & kW & 本地 PV 预测 \\
    $c_k$ & 分时电价 & CNY/kWh 或 USD/MWh & 经济调度信号 \\
    \bottomrule
  \end{tabular}
\end{table}

这些变量不由控制器直接改变，而是作为 MPC 预测窗口内的已知或预测输入。实际实验中，PV 预测误差可以通过对 $P_{PV,k}$ 添加扰动构造。

\section{控制变量与状态变量}

\begin{table}[htbp]
  \centering
  \caption{状态变量、控制变量和辅助变量}
  \begin{tabular}{llll}
    \toprule
    类型 & 符号 & 含义 & 单位 \\
    \midrule
    状态变量 & $SOC_k$ & TES 荷冷状态 & - \\
    状态变量 & $T_k$ & 机房温度 & \si{\degreeCelsius} \\
    控制变量 & $q_{ch,k}$ & TES 充冷功率 & kW \\
    控制变量 & $q_{dis,k}$ & TES 放冷功率 & kW \\
    控制变量 & $P_{grid,k}$ & 电网购电功率 & kW \\
    控制变量 & $P_{spill,k}$ & PV 弃电功率 & kW \\
    二进制变量 & $z_{ch,k}$ & 充冷模式 & - \\
    二进制变量 & $z_{dis,k}$ & 放冷模式 & - \\
    松弛变量 & $s_{T,k}$ & 温度约束松弛 & \si{\degreeCelsius} \\
    松弛变量 & $s_{SOC,k}$ & SOC 约束松弛 & - \\
    \bottomrule
  \end{tabular}
\end{table}

\section{TES 蓄冷系统模型}

令 $E_{TES}$ 表示 TES 有效蓄冷容量，$\Delta t$ 表示控制步长，$\eta_{ch}$ 和 $\eta_{dis}$ 分别表示充冷和放冷效率，$\lambda$ 表示自损耗系数。TES SOC 动态为：

\begin{equation}
  SOC_{k+1}
  =
  (1-\lambda \Delta t)SOC_k
  + \frac{\eta_{ch} q_{ch,k}\Delta t}{E_{TES}}
  - \frac{q_{dis,k}\Delta t}{\eta_{dis}E_{TES}}.
  \label{eq:soc}
\end{equation}

SOC 需要满足安全范围：

\begin{equation}
  SOC_{\min} - s_{SOC,k} \leq SOC_k \leq SOC_{\max} + s_{SOC,k}, \quad s_{SOC,k}\geq 0.
\end{equation}

充放冷功率限制为：

\begin{align}
  0 \leq q_{ch,k} &\leq q_{ch}^{max}, \\
  0 \leq q_{dis,k} &\leq q_{dis}^{max}.
\end{align}

为避免同一时刻同时充冷和放冷，引入二进制变量：

\begin{align}
  q_{ch,k} &\leq z_{ch,k}q_{ch}^{max}, \\
  q_{dis,k} &\leq z_{dis,k}q_{dis}^{max}, \\
  z_{ch,k}+z_{dis,k} &\leq 1, \\
  z_{ch,k},z_{dis,k} &\in \{0,1\}.
\end{align}

\section{机房温度代理模型}

为了使 MPC 问题能够高频滚动求解，本文采用控制导向温度代理模型描述机房温度动态：

\begin{equation}
  T_{k+1}
  =
  aT_k
  + bT_{out,k}
  + cP_{IT,k}
  - d(q_{base,k}+q_{dis,k}).
  \label{eq:temp}
\end{equation}

其中 $a$ 描述机房温度惯性，$b$ 描述室外温度影响，$c$ 描述 IT 负荷发热影响，$d$ 描述冷却量对机房温度的抑制作用，$q_{base,k}$ 表示基础冷却量。温度约束写为：

\begin{align}
  T_k &\leq T_{\max}+s_{T,k}, \\
  T_k &\geq T_{\min}-s_{T,k}, \\
  s_{T,k} &\geq 0.
\end{align}

温度松弛变量只用于避免优化问题不可行，并在目标函数中给予高权重惩罚。最终结果分析中应报告温度越界次数和越界 degree-hours。

\section{光伏-电网-冷却系统功率平衡模型}

系统电功率平衡可写为：

\begin{equation}
  P_{grid,k} - P_{spill,k}
  =
  P_{IT,k}
  + P_{cold,k}
  - P_{PV,k}.
  \label{eq:power_balance}
\end{equation}

其中 $P_{IT,k}$ 为 IT 负荷，$P_{cold,k}$ 为冷站电功率，$P_{PV,k}$ 为 behind-the-meter PV 出力。TES 充放冷通过改变冷机出力和冷站电功率间接影响 $P_{cold,k}$。电网购电和 PV 弃电满足：

\begin{equation}
  P_{grid,k}\geq 0,\quad P_{spill,k}\geq 0.
\end{equation}

该模型采用全设施侧 PV/grid 口径。代码中仍保留冷站侧 proxy grid、proxy spill 和 proxy cost 指标，用于和旧结果对照及解释冷站控制贡献，但论文主成本、主购电量和主 PV spill 均采用全设施侧口径。

\section{运行成本与 PUE 计算模型}

运行成本定义为：

\begin{equation}
  J_{cost}=\sum_{k} c_k P_{grid,k}\Delta t.
\end{equation}

PUE 可按总设施功率和 IT 功率计算：

\begin{equation}
  PUE_k=\frac{P_{IT,k}+P_{cooling,k}+P_{aux,k}}{P_{IT,k}}.
\end{equation}

最终实验结果不应只报告总成本，还应报告 PV 弃电、PV 自消纳率、PUE、温度越界、SOC 行为和求解器性能。

\section{模型假设与参数设置}

\begin{table}[htbp]
  \centering
  \caption{主要参数设置草案}
  \begin{tabular}{lll}
    \toprule
    参数 & 数值或目标设置 & 含义 \\
    \midrule
    $\Delta t$ & 0.25 h & 15 min 控制步长 \\
    $N$ & 48 & 12 h 预测窗口 \\
    $E_{TES}$ & 18000 kWh & TES 有效蓄冷容量 \\
    $q_{ch}^{max}$ & 4500 kW & 最大充冷功率 \\
    $q_{dis}^{max}$ & 4500 kW & 最大放冷功率 \\
    $T_{\min},T_{\max}$ & 23.5--27 \si{\degreeCelsius} & 机房温度约束 \\
    Solver & HiGHS & MILP 求解器目标 \\
    \bottomrule
  \end{tabular}
\end{table}

上述参数已与本地 \texttt{mpc\_v2/config/base.yaml} 的默认配置对齐。当前本地可验证代码包括南京 EnergyPlus 最小模型包，以及顶层 \texttt{mpc\_v2} 中的确定性 MILP-MPC 闭环验证包。

\chapter{基于 MILP-MPC 的控制方法设计}

\section{MPC 控制思想}

MPC 在每个控制时刻执行以下步骤：

\begin{enumerate}[label=(\arabic*)]
  \item 获取当前状态 $SOC_k$ 和 $T_k$。
  \item 生成未来 $N$ 步的 PV、电价、室外温度和 IT 负荷预测。
  \item 求解未来 $N$ 步的 MILP 优化问题。
  \item 只执行第一步控制动作。
  \item 更新系统状态并进入下一控制时刻。
\end{enumerate}

这种滚动优化机制使控制器可以随着预测更新不断修正策略。即使远期 PV 或天气预测存在误差，下一时刻的重新求解也可以部分修正前一步决策。

\section{预测时域与滚动优化机制}

本文当前候选实验采用 15 min 控制步长和 12 h 预测时域，对应 $N=48$。该设置能够覆盖半日尺度电价和 PV 变化，并使 7 day 闭环归因实验可以在本地稳定复现。48 h 预测时域对应 $N=192$，仅作为 slow/manual 扩展实验保留，不作为当前最终候选结果口径。

实际实现时，应记录每次求解的状态、预测输入、第一步动作、目标函数值、求解状态和求解时间，保证结果可复现。

\section{优化变量定义}

在预测窗口 $\mathcal{K}=\{0,\ldots,N-1\}$ 内，主要优化变量包括：

\begin{equation}
  \mathbf{u}
  =
  \{q_{ch,k},q_{dis,k},P_{grid,k},P_{spill,k},z_{ch,k},z_{dis,k},s_{T,k},s_{SOC,k}\}_{k=0}^{N-1}.
\end{equation}

状态变量 $SOC_k$ 和 $T_k$ 通过动态方程在预测窗口内递推。二进制变量 $z_{ch,k}$ 和 $z_{dis,k}$ 使问题成为 MILP。

\section{目标函数设计}

MPC 目标函数可写为：

\begin{align}
  \min \sum_{k=0}^{N-1}
  \Big[
    c_k P_{grid,k}\Delta t
    + \alpha P_{spill,k}\Delta t
    + \beta(q_{ch,k}+q_{dis,k})\Delta t
    + \rho_T s_{T,k}
    + \rho_{SOC}s_{SOC,k}
    + \rho_{\Delta u} r_{\Delta u,k}
  \Big]
  + \rho_f s_f.
  \label{eq:objective}
\end{align}

其中第一项最小化购电成本；第二项惩罚 PV 弃电；第三项避免 TES 过度循环；第四、五项惩罚温度和 SOC 约束松弛；第六项减少控制动作频繁切换；最后一项用于约束终端 SOC 或终端偏差。

\begin{table}[htbp]
  \centering
  \caption{目标函数各项解释}
  \begin{tabular}{lll}
    \toprule
    项 & 作用 & 说明 \\
    \midrule
    $c_kP_{grid,k}\Delta t$ & 购电成本 & 经济性主目标 \\
    $\alpha P_{spill,k}\Delta t$ & PV 弃电惩罚 & 提高 PV 消纳 \\
    $\beta(q_{ch,k}+q_{dis,k})\Delta t$ & TES 循环惩罚 & 避免无意义充放冷 \\
    $\rho_T s_{T,k}$ & 温度松弛惩罚 & 保证机房安全 \\
    $\rho_{SOC}s_{SOC,k}$ & SOC 松弛惩罚 & 保证 TES 安全范围 \\
    $\rho_{\Delta u} r_{\Delta u,k}$ & 动作变化惩罚 & 减少频繁切换 \\
    $\rho_f s_f$ & 终端惩罚 & 避免预测窗口末端透支 TES \\
    \bottomrule
  \end{tabular}
\end{table}

\section{系统约束条件}

\begin{longtable}{lll}
  \caption{MPC 约束条件汇总} \\
  \toprule
  约束类型 & 数学形式 & 物理意义 \\
  \midrule
  \endfirsthead
  \toprule
  约束类型 & 数学形式 & 物理意义 \\
  \midrule
  \endhead
  SOC 动态 & 式 \eqref{eq:soc} & 描述 TES 能量变化 \\
  SOC 范围 & $SOC_{\min}\leq SOC_k\leq SOC_{\max}$ & TES 安全运行范围 \\
  充冷限制 & $0\leq q_{ch,k}\leq q_{ch}^{max}$ & 设备能力限制 \\
  放冷限制 & $0\leq q_{dis,k}\leq q_{dis}^{max}$ & 设备能力限制 \\
  充放互斥 & $z_{ch,k}+z_{dis,k}\leq 1$ & 避免同时充放冷 \\
  温度动态 & 式 \eqref{eq:temp} & 描述机房热动态 \\
  温度范围 & $T_{\min}\leq T_k\leq T_{\max}$ & 保证机房安全 \\
  功率平衡 & 式 \eqref{eq:power_balance} & 保证电功率平衡 \\
  非负约束 & $P_{grid,k},P_{spill,k}\geq 0$ & 物理可行性 \\
  \bottomrule
\end{longtable}

\section{混合整数线性规划表述}

由于 TES 不能同时充冷和放冷，本文引入二进制变量表示充放冷模式。结合线性 SOC 动态、线性温度代理模型、线性功率平衡和分段可线性化的动作变化惩罚，整个优化问题可以表述为 MILP。该形式便于使用 HiGHS 等开源求解器，并能够提供明确的求解状态，例如 optimal、time limit 或 infeasible。

\section{闭环控制流程}

闭环仿真应保存以下信息：

\begin{enumerate}[label=(\arabic*)]
  \item 当前时间戳、状态和预测输入。
  \item MILP 求解状态、目标函数值和求解时间。
  \item 第一时刻控制动作。
  \item 更新后的 SOC、温度、购电功率和 PV 弃电。
  \item 每个场景的汇总指标。
\end{enumerate}

这样可以保证论文结果可追溯，也便于定位某一场景下的异常行为。

\section{求解器与软件实现}

控制层已在顶层 \texttt{mpc\_v2} 目录中使用 Python 实现，MILP 求解器为 SciPy/HiGHS。输入包括 PV、TOU 电价、室外温度和负荷时间序列；输出包括逐时刻控制动作、状态轨迹、求解日志和指标汇总。当前仓库同时包含三类验证路径：合成/回放闭环、EnergyPlus-derived 年度 profile replay，以及 EnergyPlus Runtime API 在线耦合 pilot。论文中的最终容量矩阵应以后续完整在线 MPC+EnergyPlus 年度结果为准。

2026 年 5 月 6 日之后的本地实现进一步加入了中国分时电价/尖峰电价情景服务和 DR/peak-cap 事件服务。电价服务支持将总电价拆分为可浮动能量价和非浮动附加价，并通过峰谷价差系数、尖峰上浮率和浮动比例构造政策锚定的敏感性实验。DR 服务支持事件窗口、baseline、请求削减量、动态购电 cap、响应率和事件收益估算。上述扩展用于后续实验设计和日志记录；在完整矩阵运行前，其结果仍不能写成统计显著结论。

\chapter{实验设计与评价指标}

\section{实验平台与代码结构}

当前可验证模型位于 \texttt{Nanjing-DataCenter-TES-EnergyPlus}，其包含南京 EnergyPlus 模型、EPW 天气、TOU 电价和 PV 输入。确定性 MPC 控制层位于仓库顶层 \texttt{mpc\_v2}，采用如下结构：

\begin{verbatim}
mpc_v2/
  config/
    base.yaml
    scenario_sets.yaml
  core/
    tes_model.py
    room_model.py
    facility_model.py
    tariff_service.py
    dr_service.py
    mpc_problem_milp.py
    controller.py
    forecast.py
    metrics.py
    statistics.py
  scripts/
    run_closed_loop.py
    run_validation_matrix.py
    analyze_results.py
\end{verbatim}

配套测试位于仓库顶层 \texttt{tests/}。该控制层复用 \texttt{Nanjing-DataCenter-TES-EnergyPlus/inputs} 中的 PV、PVGIS 和 TOU CSV，不恢复旧 RL 训练目录；在线耦合代码位于 \texttt{Nanjing-DataCenter-TES-EnergyPlus/controller/energyplus_mpc/} 和 \texttt{mpc\_v2/scripts/run\_phase3\_online\_mpc\_ep\_matrix.py}。

\section{输入数据与参数设置}

输入数据包括南京 EPW 天气、江苏/南京 TOU 电价、南京 PV 预测和 IT 负荷假设。当前 EnergyPlus 模型已验证 15 min timestep；控制层也应使用 15 min 步长以保持一致。

\begin{table}[htbp]
  \centering
  \caption{实验参数草案}
  \begin{tabular}{lll}
    \toprule
    参数 & 数值 & 含义 \\
    \midrule
    控制步长 & 0.25 h & 15 min \\
    预测时域 & 12 h & MPC horizon \\
    预测步数 & 48 & 每次优化步数 \\
    TES 容量 & 18000 kWh & 有效蓄冷容量 \\
    最大充冷功率 & 4500 kW & TES charging limit \\
    最大放冷功率 & 4500 kW & TES discharging limit \\
    温度范围 & 23.5--27 \si{\degreeCelsius} & 机房约束 \\
    \bottomrule
  \end{tabular}
\end{table}

\section{Baseline 控制策略}

本文设置 direct no-TES、price-aware no-TES MPC、rule-based TES 和 TES-MPC 四类 baseline / controller，以拆分冷机经济调度和 TES 的增量贡献：

\begin{table}[htbp]
  \centering
  \caption{控制策略对比}
  \begin{tabular}{lll}
    \toprule
    控制方法 & 说明 & 用途 \\
    \midrule
    direct no-TES & 不使用蓄冷，按温度需求直接调冷机 & 基础对照 \\
    MPC no-TES & 使用同一 MILP 和价格/PV/grid 目标，但禁用 TES & 归因冷机经济调度 \\
    rule-based TES & 按固定低价/高价时段充放冷 & 简单规则对照 \\
    TES-MPC & 根据预测信息滚动优化冷机与 TES & 本文方法 \\
    \bottomrule
  \end{tabular}
\end{table}

MPC no-TES 是当前最关键的科学归因 baseline。只有 TES-MPC 相对 MPC no-TES 的差异，才可解释为 TES 的增量贡献。

\section{MILP-MPC 控制器设置}

MPC 控制器每 15 min 求解一次 12 h MILP 问题，只执行第一步动作。求解器状态和求解时间必须记录。若求解失败，应采用保守 fallback，例如停止 TES 放冷并满足温度安全约束，同时记录 infeasible 或 solver failure 次数。

\section{场景矩阵设计}

\begin{table}[htbp]
  \centering
  \caption{实验场景矩阵}
  \begin{tabular}{lll}
    \toprule
    场景 & 目的 & 结果关注点 \\
    \midrule
    nominal & 基本控制行为 & SOC、充放冷、温度、PV/grid \\
    baseline\_no\_tes & 无 TES 对照 & 成本和 PUE 基线 \\
    rule\_based\_tes & 规则控制对照 & MPC 增量 \\
    pv\_error\_5 & 轻微预测误差 & 鲁棒性 \\
    pv\_error\_10 & 中等预测误差 & PV spill 和成本变化 \\
    pv\_error\_20 & 较强预测误差 & 鲁棒性边界 \\
    hot\_week & 高温天气 & 温度约束和 TES 透支风险 \\
    mild\_week & 温和天气 & TES 利用率 \\
    tariff\_low & 低电价价差 & TES 经济性下限 \\
    tariff\_base & 基准电价 & 标准对比 \\
    tariff\_high & 高电价价差 & TES 经济性上限 \\
    \bottomrule
  \end{tabular}
\end{table}

\section{评价指标}

\begin{table}[htbp]
  \centering
  \caption{评价指标定义}
  \begin{tabular}{lll}
    \toprule
    类别 & 指标 & 说明 \\
    \midrule
    经济性 & total electricity cost & 总购电成本 \\
    经济性 & cost saving ratio & 相对 baseline 节省率 \\
    经济性 & grid import & 电网购电量 \\
    经济性 & peak grid demand & 峰值购电功率 \\
    PV 利用 & PV self-consumption ratio & PV 自消纳比例 \\
    PV 利用 & PV load coverage ratio & PV 对负荷覆盖比例 \\
    PV 利用 & PV spill & PV 弃电量 \\
    温度安全 & violation count & 温度越界次数 \\
    温度安全 & violation degree-hours & 温度越界程度 \\
    TES 行为 & equivalent cycles & TES 等效循环次数 \\
    求解性能 & average solve time & 平均求解时间 \\
    求解性能 & optimal solve rate & 最优求解率 \\
    \bottomrule
  \end{tabular}
\end{table}

\section{可复现性说明}

每次实验应保存配置文件、输入数据版本、随机种子、求解器版本、运行命令和结果摘要。论文中的所有表格和图都应能从保存的结果文件重新生成。对于尚未完成的实验，正文只能保留占位，不应写入未经验证的数值结论。

\chapter{结果分析与讨论}

本章记录当前 \texttt{chiller\_tes\_mpc\_attribution\_20260505} 版本的可复现闭环仿真结果。结果文件保存于 \texttt{results/chiller\_tes\_mpc\_attribution\_20260505/}，对应核心归因矩阵均采用 15 min 控制步长、672 个闭环步和 48 步滚动预测时域，即 7 day 闭环、12 h prediction horizon。新增中国 TOU/DR 1 个月矩阵结果保存于 \texttt{results/china\_tou\_dr\_matrices\_20260506/}，共 138 个 30 day run，每个 run 为 2880 个 15 min 闭环步。当前主成本和 PV/grid 指标采用全设施 behind-the-meter 口径；冷站侧 proxy 指标仅用于解释冷站控制贡献和与历史结果对照。

\section{Nominal case 下的 MPC 调度行为}

Nominal case 首先用于检查 MPC 是否产生合理行为，而不是直接证明全部经济收益均来自 TES。本版本中，MPC 能够调用冷机 mode、维持机房温度约束，并形成 TES 双向充放冷。7 day TES-MPC 中，TES 充冷约 21676.87 kWh\_th，放冷约 22821.69 kWh\_th，充冷加权电价约 29.60 USD/MWh，放冷加权电价约 186.80 USD/MWh，说明存在低价充冷、高价放冷行为。

分析时需要回答：

\begin{enumerate}[label=(\arabic*)]
  \item 控制器是否在低价或高 PV 时段倾向于充冷？
  \item 控制器是否在高价时段或 PV 不足时段倾向于放冷？
  \item SOC 是否保持在规划范围内？
  \item 机房温度是否满足约束？
  \item 控制动作是否出现过度频繁切换？
\end{enumerate}

对当前结果的回答是：控制器会在低价时段充冷、在高价时段放冷，SOC 保持在物理边界内，温度约束满足；但 inventory-using TES-MPC 最终 SOC 从 0.50 降至 0.15，因此该结果必须报告库存消耗，不能假装为周期中性结果。

\section{归因 baseline 对比}

\begin{table}[htbp]
  \centering
  \caption{7 day 归因矩阵主要结果}
  \begin{tabular}{llllll}
    \toprule
    指标 & direct no-TES & MPC no-TES & RBC TES & MPC TES & SOC-neutral attempt \\
    \midrule
    Total cost & 303552.69 & 289706.52 & 303263.15 & 289447.59 & 289400.14 \\
    Grid import (kWh) & 3251529.47 & 3239731.92 & 3255772.89 & 3238814.27 & 3238582.23 \\
    Peak grid (kW) & 20633.25 & 21090.00 & 20576.55 & 21090.00 & 21090.00 \\
    PV spill (kWh) & 0.00 & 0.00 & 0.00 & 0.00 & 0.00 \\
    Cold-station energy (kWh) & 336516.59 & 324719.04 & 340760.01 & 323801.39 & 323569.35 \\
    TES charge/discharge (kWh\_th) & 0/0 & 0/0 & 37012.50/27562.50 & 21676.87/22821.69 & 43477.06/41014.60 \\
    Final SOC & 0.50 & 0.50 & 0.56 & 0.15 & 0.15 \\
    \bottomrule
  \end{tabular}
\end{table}

该对比显示，price-aware no-TES MPC 相对 direct no-TES 的总成本下降约 4.56\%，说明冷机经济调度、热惯性利用和 mode 优化已经解释了主要经济收益。TES-MPC 相对 no-TES MPC 的增量总成本下降约 0.09\%，这才是当前 nominal 归因矩阵中 TES 的增量经济贡献。RBC TES 能够形成明显充放冷和较低峰值，但总成本接近 direct no-TES。SOC-neutral attempt 未能实现库存中性，最终 SOC 仍为 0.15，因此只能作为失败诊断结果。

\section{与 rule-based TES 的对比}

RBC 在本版本中产生了明显 TES 充放冷：7 day 充冷约 37012.50 kWh\_th、放冷约 27562.50 kWh\_th，且最终 SOC 略高于初始值。RBC 的峰值低于 direct no-TES 和 MPC，但总成本只比 direct no-TES 略低。MPC TES 的总成本最低，但峰值高于 direct no-TES，说明 energy-charge-only 目标下 MPC 倾向于经济调度而非自然削峰；削峰结论必须放在 peak-cap 或 demand-charge 场景中讨论。

\section{PV 预测误差鲁棒性分析}

PV error 场景用于检验 MPC 对可再生能源预测误差的敏感性。当前归因矩阵改为全设施 behind-the-meter PV/grid 口径后，nominal 6 MWp PV 相对于 18 MW IT 负荷较小，因此 7 day attribution cases 的主 PV spill 均为 0。这说明 nominal attribution matrix 不能用于证明 TES 提高 PV 消纳。旧冷站侧 proxy PV spill 仍保存在结果文件中，但只作为辅助诊断，不作为论文主 PV 结论。

\section{气象条件敏感性分析}

hot week 和 mild week 场景用于分析天气条件对 TES-MPC 的影响。当前归因修复轮未重新扩大天气矩阵，因此本节不再引用旧 24 h smoke 数字作为最终结论。后续若进入论文定稿，应使用与 attribution matrix 相同的全设施 PV/grid 口径和 12 h horizon 重新生成 hot/mild week。

\section{分时电价敏感性分析}

TES 经济性高度依赖电价价差。当前 7 day attribution matrix 中，TES-MPC 的充冷加权电价低于放冷加权电价，说明控制器存在电价套利行为。但 TES 的增量降本只有约 0.09\%，因此不能把 direct no-TES 到 TES-MPC 的全部降本归因于 TES。需求电费或 peak-cap 场景需要单独报告，不能与 energy-charge-only nominal case 混写。

在 30 day 中国 TOU screening 矩阵中，\texttt{gamma=\{0,0.5,1.0,1.5,2.0\}}、\texttt{cp\_uplift=\{0,0.2\}}、hot/mild weather 和 \texttt{mpc\_no\_tes}/\texttt{mpc} 两类控制器形成 40 个场景。该矩阵中 TES-MPC 平均月成本为 1,291,919.83 CNY，no-TES MPC 平均月成本为 1,292,366.36 CNY；配对平均节省为 446.52 CNY。完整 TOU 四控制器对比中，TES-MPC 平均月成本为 1,296,598.71 CNY，no-TES MPC 为 1,297,248.94 CNY，rule-based TES 为 1,347,132.29 CNY，direct no-TES 为 1,348,448.17 CNY。该结果支持“TES 有可测增量收益，但量级小于价格感知 MPC 本身收益”的表述。

\section{PV--TES 技术容量矩阵}

为避免把单一 18 MWh\_th TES 参考值写成最优容量，本文新增三地点 PV--TES 技术容量矩阵。当前已完成的 150 个全年场景使用南京、广州和北京真实 EPW 驱动的 EnergyPlus 年度 no-control 负荷/天气边界，PV 出力来自 PVGIS v5.3 20 MWp 曲线并按装机线性缩放，电价统一采用江苏 2025 TOU 曲线。主扫描为 $PV=\{0,10,20,40,60\}$ MWp、$TES=\{0,9,18,36,72\}$ MWh\_th、尖峰电价上浮 $\delta=\{0,0.2\}$，每个场景包含 35040 个 15 min 控制步。但该版本是 EnergyPlus-derived profile replay：EnergyPlus 提供年度负荷/天气边界，MPC-style TES dispatch 在该边界上回放，不是 MPC 动作实时写入 EnergyPlus 后得到 plant response 的在线耦合矩阵。

审计结果显示 EnergyPlus-derived 全量矩阵 150 个场景均完成，P0 错误为 0；唯一 P1 警告是 15 个 9 MWh\_th 小 TES 场景出现轻微负尖峰电价抑制，说明容量过小且回充惩罚可能超过尖峰放冷收益。三地点最大尖峰电价抑制率分别约为北京 0.587、广州 0.528、南京 0.482；20 MWp PV + 72 MWh\_th TES 在该筛选边界下分别取得北京 0.439、广州 0.365、南京 0.352 的尖峰抑制率，同时 PV 自消纳率分别约为 0.829、0.947、0.912。年度绝对峰值削减接近零，说明当前 EnergyPlus-derived 边界下主收益并非年度峰值迁移，而是尖峰窗口购电抑制和 PV 自消纳改善。

为修正上述边界，当前代码已新增在线 MPC+EnergyPlus runner。在线路径中，MPC 每 15 min 通过 EnergyPlus Runtime API 写入 TES 控制动作，EnergyPlus 返回机房温度、蓄冷罐换热、冷机功率和全设施电功率。南京 pilot、北京 smoke 和广州 smoke 已验证在线链路可运行且温度越界比例为 0；南京 TES=9 MWh\_th、尖峰上浮 20\% 的 pilot 中，TES 在尖峰窗口产生 1077.75 kWh\_th 放冷，CP 电网电量相对 no-TES 降低 47.75 kWh，signed-valve 约束满足。完整三地点全年在线 MPC+EnergyPlus 容量矩阵尚未完成，因此当前论文不能把 EnergyPlus-derived 150 场景结论写作最终在线耦合结论。

因此，本节当前只能给出 EnergyPlus-derived screening 下的运行层面 technical recommended capacity range：20 MWp PV 可作为三地点共同参考装机，TES 容量宜落在 36--72 MWh\_th 的高容量区间，其中 72 MWh\_th 在该筛选规则下更接近最大尖峰抑制效果。该结论不包含 CAPEX、LCOE、NPV 或投资回收约束，不能表述为 economic optimum；在线 MPC+EnergyPlus 年度矩阵完成后，应以在线结果替换或校准本段容量推荐。

\section{DR 与 peak-cap 场景分析}

DR event 矩阵包含 day-ahead、fast 和 realtime 三类 notice，每类分别测试 5\%、10\% 和 15\% 请求削减率，并对比 no-TES MPC 与 TES-MPC。30 day episode 中每个场景只触发一次 DR 事件。当前代理模型下，请求削减量分别为 1900、3800 和 5700 kWh，达成削减量为 2060.78 kWh，对应响应率为 1.0846、0.5423 和 0.3615。由于本版本尚未核验真实需求响应 baseline 和补偿结算规则，DR 收益只能表述为情景估算。

Peak-cap 矩阵使用 no-cap no-TES MPC baseline 自动推导峰值，再按 \texttt{r\_cap=\{1.00,0.99,0.97,0.95\}} 缩放。正式结果中对应达成峰值约为 21090.0、20879.1、20457.3 和 20035.5 kW，\texttt{peak\_slack\_max\_kw} 近似为 0，说明约束在优化层面被满足。但最严格的 \texttt{r\_cap=0.95} 场景出现约 415--431 degree-hours 温度违约，应作为强削峰压力测试讨论，不能作为推荐控制点。

\section{温度约束与 SOC 约束分析}

温度和 SOC 约束决定控制策略是否工程可接受。7 day attribution matrix 中，direct no-TES、MPC no-TES、RBC TES、MPC TES 和 SOC-neutral attempt 均满足 fallback count 为 0、温度越界 degree-hours 为 0、物理一致性违约为 0。SOC 方面，MPC TES 会从 0.50 降至 0.15，说明其使用了初始库存；SOC-neutral attempt 在当前滚动终端惩罚下仍未实现库存中性，因此后续若要得到周期中性结论，需要加入 episode-end terminal constraint 或末端 horizon 截断策略。

\section{求解性能分析}

MPC 需要在每个控制步反复求解 MILP，因此求解性能是工程可行性的重要指标。当前默认测试不再使用 192-step horizon，而是采用 48-step horizon，使 \texttt{pytest -q} 能快速完成。7 day attribution matrix 采用 12 h receding horizon；48 h prediction horizon 只保留为 slow/manual 扩展实验，不能写成最终候选 horizon。

\section{经济收益受限原因与适用边界讨论}

如果实验显示成本节省有限，可能原因包括：

\begin{enumerate}[label=(\arabic*)]
  \item PV 规模相对数据中心负荷过小，能够转移的可再生电力有限。
  \item TOU 电价价差不足，无法覆盖 TES 充放冷效率损失。
  \item TES 容量或功率与冷却负荷不匹配。
  \item 温度约束和终端 SOC 约束限制了自由调度空间。
  \item 冷机 COP 或温度代理模型过于简化。
\end{enumerate}

本版本最重要的边界是：MPC nominal case 的大部分经济收益来自价格感知冷机调度、热惯性利用和冷机 mode 优化，而不是 TES 的单独贡献。TES 在 7 day MPC 中确实双向充放冷，但相对 no-TES MPC 的增量降本较小。后续需要更严格的 episode-end SOC 约束、需求峰值结算窗口或针对 PV surplus 的场景设计，才能证明 TES 的库存中性削峰填谷价值。

\chapter{结论与展望}

\section{主要研究结论}

本文建立了面向数据中心冷却系统的冷机 + TES-PV-TOU 控制导向模型，并构建了基于 MILP 的 MPC 控制器。该控制器能够在滚动预测时域内优化冷机出力、TES 充冷、购电功率和 PV 弃电，同时考虑机房温度约束与 TES SOC 约束。

当前闭环结果表明，冷机 + TES MPC 可以在 7 day attribution matrix 中稳定运行，并保持温度约束、安全约束和 signed-valve 约束。新增 no-TES MPC baseline 后可以看出，direct no-TES 到 TES-MPC 的大部分降本来自价格感知冷机调度和 mode 优化；TES 相对 no-TES MPC 的增量降本约 0.09\%。在 30 day 中国 TOU/DR 正式矩阵中，138 个 run 全部完成，最低 feasible rate 为 1.0，仅出现 1 次 fallback；61 组配对场景中 TES-MPC 相对 no-TES MPC 的平均月节省为 518.81 CNY。三地点 EnergyPlus-derived PV--TES 技术容量矩阵进一步给出 20 MWp PV 与 36--72 MWh\_th TES 的筛选参考区间，并显示当前边界下尖峰窗口抑制和 PV 自消纳改善比年度绝对削峰更明显。当前新增的在线 MPC+EnergyPlus pilot 已验证 Runtime API 耦合链路、温度安全和 TES 物理响应，但完整三地点年度在线矩阵尚未完成。因此，当前版本已经达到“模型清楚、变量清楚、闭环可运行、baseline 可比较、归因可拆分、在线耦合链路可验证”的阶段目标，但尚未达到“完成三地点全年在线 MPC+EnergyPlus sizing 矩阵并证明 TES 在库存中性条件下显著削峰填谷或显著提高 PV 消纳”的最终目标。

\section{本文不足}

当前研究仍存在以下不足：

\begin{enumerate}[label=(\arabic*)]
  \item 控制导向温度模型相比完整 EnergyPlus 模型更简化，物理细节有限。
  \item 冷机 COP、部分负荷效率和辅助设备能耗模型仍需进一步校准。
  \item PV、天气和负荷预测误差目前采用场景扰动方式，尚未建立概率预测模型。
  \item 当前实验仍以仿真为主，缺少真实数据中心运行数据验证。
\end{enumerate}

\section{后续工作展望}

后续可从以下方向扩展：

\begin{enumerate}[label=(\arabic*)]
  \item 将 MPC 控制层与 EnergyPlus 模型闭环耦合，提升物理保真度。
  \item 引入更真实的冷机 COP 和部分负荷模型。
  \item 使用 stochastic MPC 或 robust MPC 处理预测不确定性。
  \item 联合考虑电池、动态碳强度信号和需求响应。
  \item 在更长时间尺度和更多天气年份上评估 TES-MPC 的适用边界。
\end{enumerate}

\bibliographystyle{plainnat}
\bibliography{references}

\appendix

\chapter{主要模型参数}

\begin{table}[htbp]
  \centering
  \caption{主要模型参数汇总}
  \begin{tabular}{lll}
    \toprule
    参数 & 数值或占位 & 说明 \\
    \midrule
    $\Delta t$ & 0.25 h & 控制步长 \\
    $N$ & 48 & 预测步数 \\
    $E_{TES}$ & 18000 kWh & TES 容量 \\
    $q_{ch}^{max}$ & 4500 kW & 最大充冷功率 \\
    $q_{dis}^{max}$ & 4500 kW & 最大放冷功率 \\
    $\eta_{ch}$ & 待配置填入 & 充冷效率 \\
    $\eta_{dis}$ & 待配置填入 & 放冷效率 \\
    $\lambda$ & 待配置填入 & TES 自损耗系数 \\
    $T_{\min}$ & 18 \si{\degreeCelsius} & 温度下限 \\
    $T_{\max}$ & 27 \si{\degreeCelsius} & 温度上限 \\
    \bottomrule
  \end{tabular}
\end{table}

\chapter{MPC 优化问题完整形式}

MPC 完整形式可概括为：

\begin{align}
  \min_{\mathbf{u}} \quad
  & \sum_{k=0}^{N-1}
  \Big[
    c_k P_{grid,k}\Delta t
    + \alpha P_{spill,k}\Delta t
    + \beta(q_{ch,k}+q_{dis,k})\Delta t
    + \rho_T s_{T,k}
    + \rho_{SOC}s_{SOC,k}
    + \rho_{\Delta u} r_{\Delta u,k}
  \Big] + \rho_f s_f \\
  \text{s.t.}\quad
  & SOC_{k+1}
  =
  (1-\lambda \Delta t)SOC_k
  + \frac{\eta_{ch} q_{ch,k}\Delta t}{E_{TES}}
  - \frac{q_{dis,k}\Delta t}{\eta_{dis}E_{TES}}, \\
  & T_{k+1}
  =
  aT_k
  + bT_{out,k}
  + cP_{IT,k}
  - d(q_{base,k}+q_{dis,k}), \\
  & P_{grid,k} - P_{spill,k}
  =
  P_{IT,k}
  + P_{cold,k}
  - P_{PV,k}, \\
  & SOC_{\min}-s_{SOC,k} \leq SOC_k \leq SOC_{\max}+s_{SOC,k}, \\
  & T_{\min}-s_{T,k} \leq T_k \leq T_{\max}+s_{T,k}, \\
  & 0 \leq q_{ch,k} \leq z_{ch,k}q_{ch}^{max}, \\
  & 0 \leq q_{dis,k} \leq z_{dis,k}q_{dis}^{max}, \\
  & z_{ch,k}+z_{dis,k}\leq 1, \\
  & P_{grid,k}\geq 0,\quad P_{spill,k}\geq 0, \\
  & s_{T,k}\geq 0,\quad s_{SOC,k}\geq 0, \\
  & z_{ch,k},z_{dis,k}\in\{0,1\}.
\end{align}

\chapter{实验场景配置}

本附录用于记录最终 \texttt{scenario\_sets.yaml} 或等价配置。当前占位如下：

\begin{verbatim}
nominal:
  description: base weather, base PV, base tariff
baseline_no_tes:
  tes_enabled: false
rule_based_tes:
  controller: rule_based
pv_error_5:
  pv_forecast_error: 0.05
pv_error_10:
  pv_forecast_error: 0.10
pv_error_20:
  pv_forecast_error: 0.20
hot_week:
  weather_slice: hot_week
mild_week:
  weather_slice: mild_week
tariff_low:
  tariff_spread: low
tariff_base:
  tariff_spread: base
tariff_high:
  tariff_spread: high
\end{verbatim}

\chapter{代码结构与运行命令}

当前 EnergyPlus 最小模型运行命令示例：

\begin{verbatim}
cd Nanjing-DataCenter-TES-EnergyPlus
.\run_energyplus_nanjing.ps1 -EnergyPlusExe "C:\Path\To\energyplus.exe"
\end{verbatim}

后续控制层运行命令建议：

\begin{verbatim}
python -m controller.scripts.run_closed_loop --config config/base.yaml
python -m controller.scripts.run_scenario_matrix --config config/scenario_sets.yaml
python -m controller.scripts.make_plots --results results/
python -m controller.scripts.make_report_tables --results results/
\end{verbatim}

\chapter{补充实验图表清单}

最终论文建议至少准备：

\begin{enumerate}[label=(\arabic*)]
  \item 数据中心冷却-TES-PV-电网系统结构图。
  \item MPC 滚动优化流程图。
  \item nominal case 下 SOC、充冷、放冷曲线。
  \item nominal case 下 PV、grid import、PV spill 曲线。
  \item 机房温度轨迹与温度约束。
  \item no-TES、rule-based TES、MPC 成本对比。
  \item PV 预测误差对成本、PV spill、温度 violation 的影响。
  \item hot week 与 mild week 结果对比。
  \item 不同电价价差下 TES-MPC 经济性对比。
\end{enumerate}

\chapter{旧技术路线说明}

项目前期曾考虑基于强化学习的数据中心冷却控制路线。但数据中心冷却系统具有强安全约束、真实系统试错成本高、训练稳定性要求高等特点。结合毕业设计周期和可验证性要求，本文将主线调整为基于模型的 MILP-MPC 控制方法。旧 DRL、Sinergym 和失败实验不进入正文主线，仅作为方法选择背景保留。

\end{document}
